\section{An entropy inequality from irreducible invariance}
\label{sec:entropy}

We use finite Shannon entropy with natural logarithms.  For a random
variable $X$ with finite support, write
\[
  \HH(X)=-\sum_x \Prob(X=x)\log\Prob(X=x),
\]
where $0\log 0=0$, and define conditional entropy by averaging over the
conditioning variable.  The chain rule is
$\HH(X,Y)=\HH(Y)+\HH(X\mid Y)$.
We will use nonnegativity of conditional mutual information.  To recall
its elementary basis, for probability vectors $p$ and $q$ the relative
entropy
\[
  \operatorname{KL}(p\Vert q)=\sum_{x:p_x>0}p_x\log\frac{p_x}{q_x}
\]
is nonnegative, with value $+\infty$ if some $p_x>0$ has $q_x=0$.
Otherwise, the inequality $-\log u\geq 1-u$ gives
\[
  \operatorname{KL}(p\Vert q)
  \geq \sum_{x:p_x>0}(p_x-q_x)\geq 0.
\]
Conditional mutual information is the average, over the conditioning
variable, of the relative entropy of a conditional joint law with
respect to the product of its marginals.  Thus
\[
  \operatorname{I}(A;B\mid C)
  =\HH(A\mid C)+\HH(B\mid C)-\HH(A,B\mid C)\geq 0.
\]
In particular, if $C$ is a deterministic function of $A$ and also of $B$,
the chain rule yields
\begin{equation}
  \operatorname{I}(A;B)
  =\HH(C)+\operatorname{I}(A;B\mid C)\geq\HH(C).
  \label{eq:common-quotient-information}
\end{equation}

Let $E$ be a finite-dimensional vector space over $\Q$, and let $X$ be
an $E$-valued random variable with finite support.  For a rational linear
subspace $V\leq E$, the notation $X\bmod V$ denotes its coset label,
which is a finite-valued random variable.  Put
\begin{equation}
  F_X(V)=\HH(X\mid X\bmod V)
        =\HH(X)-\HH(X\bmod V).
  \label{eq:quotient-entropy}
\end{equation}
The last equality holds because the coset label is a function of $X$.
In particular, $F_X(0)=0$ and $F_X(E)=\HH(X)$.

\begin{lemma}[Supermodularity]
\label{lem:supermodular}
For rational linear subspaces $U,V\leq E$,
\begin{equation}
  F_X(U)+F_X(V)
  \leq F_X(U+V)+F_X(U\cap V).
  \label{eq:entropy-supermodularity}
\end{equation}
\end{lemma}

\begin{proof}
Set $A=X\bmod U$, $B=X\bmod V$, and $C=X\bmod(U+V)$.
The pair $(A,B)$ and the label $X\bmod(U\cap V)$ determine each other
on the finite support: two support points have the same pair of labels
if and only if their difference belongs to both $U$ and $V$.
Moreover, $C$ is a deterministic function of each of $A$ and $B$.
Equation~\eqref{eq:common-quotient-information} therefore gives
\[
  \HH(X\bmod U)+\HH(X\bmod V)
  -\HH(X\bmod(U\cap V))
  \geq \HH(X\bmod(U+V)).
\]
Substituting \eqref{eq:quotient-entropy} proves the assertion.
\end{proof}

\begin{lemma}[Irreducible invariance]
\label{lem:invariant}
Suppose that $\dim E=t>0$, a group $G$ acts $\Q$-linearly and
irreducibly on $E$, and $gX$ has the same law as $X$ for every $g\in G$.
Then every rational subspace $V\leq E$ satisfies
\begin{equation}
  F_X(V)\leq \frac{\dim V}{t}\,\HH(X).
  \label{eq:invariant-entropy}
\end{equation}
\end{lemma}

\begin{proof}
A finite set has only finitely many partitions.  The partition of the
support into cosets of $V$ determines $F_X(V)$, so this quantity takes
only finitely many values as $V$ varies.  Since the positive dimensions
also belong to a finite set, the maximum
\[
  \lambda=\max_{0\neq V\leq E}\frac{F_X(V)}{\dim V}
\]
exists.  Choose a maximizing subspace $V$ of largest dimension.
Invariance of the law of $X$ gives $F_X(gV)=F_X(V)$, so $gV$ is another
maximizer for every $g\in G$.  Supermodularity and the definition of
$\lambda$ imply
\begin{align*}
  \lambda\bigl(\dim V+\dim gV\bigr)
  &\leq F_X(V+gV)+F_X(V\cap gV)\\
  &\leq\lambda\bigl(\dim(V+gV)+\dim(V\cap gV)\bigr).
\end{align*}
The bound on the intersection also holds when it is zero because
$F_X(0)=0$.  The two outer expressions coincide by the dimension formula.
Thus $V+gV$ attains the maximum slope as well.  The choice of $V$ forces
$\dim(V+gV)=\dim V$, whence $gV=V$.  Since this holds for every $g$,
irreducibility and $V\neq 0$ imply $V=E$.
Consequently $\lambda=F_X(E)/t=\HH(X)/t$, proving the bound for every
nonzero subspace; the zero subspace satisfies it directly.
The argument also applies when $\HH(X)=0$.
\end{proof}

\section{Permutation fibers and counts over the full support}
\label{sec:fibers}

Fix the positive multiplicity tuple $\mathbf d=(d_1,\ldots,d_n)$, with
$n\geq 2$ and $k=\sum_a d_a\geq 3$.  All constants in the remainder of
the proof may depend on this fixed tuple.  Define
\begin{equation}
  p_a=\frac{d_a}{k},\qquad
  M=\frac{k!}{\prod_{a=1}^n d_a!},\qquad
  h=\log M,\qquad
  \alpha_a=ka-\sum_{c=1}^n c\,d_c.
  \label{eq:support-parameters}
\end{equation}
The integers $\alpha_a$ are distinct.  Let $S\subseteq\mathbb Z^k$ be
the set of tuples containing $\alpha_a$ exactly $d_a$ times for every
$a$.  Replacing the symbol $a$ by $\alpha_a$ is only an injective
labelling of the basis symbols of $D_{\mathbf d}$; it is not an operation
on the tensor.  We have $|S|=M$ and
\[
  \sum_{i=1}^k s_i=0\qquad(s\in S).
\]
Let $Z$ be uniform on $S$.  Each coordinate $Z_j$ has distribution $p$
on the integer labels, so
\[
  \HH(Z)=h,\qquad \HH(Z_j)=\HH(p),\qquad
  \HH(p)=-\sum_{a=1}^n p_a\log p_a.
\]
For $1\leq j\leq k$, put
\begin{equation}
  E_j=\left\{v\in\Q^k:
           \sum_{i=1}^k v_i=0,\ v_j=0\right\},
  \qquad \dim E_j=k-2,
  \qquad \beta=\frac{h-\HH(p)}{k-2}\geq 0.
  \label{eq:collision-space}
\end{equation}
Nonnegativity follows from $h-\HH(p)=\HH(Z\mid Z_j)$.

\begin{lemma}[Entropy in a fixed-coordinate fiber]
\label{lem:fiber-entropy}
If $V\leq E_j$ is a rational subspace of dimension $r$, then
\begin{equation}
  F_Z(V)\leq r\beta.
  \label{eq:fiber-entropy-bound}
\end{equation}
\end{lemma}

\begin{proof}
Since every vector of $V$ has coordinate $j$ equal to zero, the quotient
label $Z\bmod V$ determines $Z_j$.  Condition on $Z_j=\alpha_a$.
The other $k-1$ coordinates are uniformly distributed among permutations
of the multiset obtained by removing one copy of $\alpha_a$; their sum
is $-\alpha_a$.  Set the $j$th coordinate to zero and replace every
other coordinate $Z_i$ by $Z_i+\alpha_a/(k-1)$.  This translates the
conditional support into $E_j$.  The translation preserves differences,
entropy, and the conditional entropy modulo $V$.  The translated law
is invariant under all permutations of the $k-1$ coordinates different
from $j$.

The resulting standard action of the symmetric group on $E_j$ is
irreducible over $\Q$.  Indeed, let $W\leq E_j$ be a nonzero invariant
subspace and choose $v\in W\setminus\{0\}$.  Among the $k-1$ moving
coordinates, two entries $v_u,v_w$ must differ, since they sum to zero.
Subtracting the image of $v$ under their transposition yields
\[
  (v_u-v_w)(e_u-e_w)\in W.
\]
The coefficient is nonzero.  Permuting this difference gives all
coordinate differences among the moving coordinates, which span $E_j$.
Thus $W=E_j$.  The same argument includes $k=3$, when $E_j$ is
one-dimensional.

Lemma~\ref{lem:invariant}, applied to each translated conditional law,
now gives
\[
  \HH(Z\mid Z\bmod V,\ Z_j=\alpha_a)
  \leq\frac{r}{k-2}\,\HH(Z\mid Z_j=\alpha_a).
\]
A conditional law supported at one point has entropy zero and obeys the
same inequality.  Average with weights
$\Prob(Z_j=\alpha_a)=p_a$.  Since $Z\bmod V$ determines $Z_j$, the left
side averages to $F_Z(V)$.  The right side averages to
\[
  \frac{r}{k-2}\,\HH(Z\mid Z_j)
  =\frac{r}{k-2}\bigl(h-\HH(p)\bigr)
  =r\beta,
\]
which proves the claim.
\end{proof}

For $s\in S$ and a rational subspace $V\leq E_j$, define
\[
  C_V(s)=(s+V)\cap S.
\]
Conditioned on a coset, the uniform variable $Z$ is uniform on the
corresponding fiber.  Therefore
\begin{equation}
  F_Z(V)=\frac{1}{M}\sum_{s\in S}\log |C_V(s)|.
  \label{eq:fiber-entropy-identity}
\end{equation}

Let $N$ be a positive multiple of $M$.  Identify a support edge of
$D_{\mathbf d}^{\boxtimes N}$ with its sequence of $N$ support columns,
and define
\[
  \mathcal A_N=
  \left\{I=(I(1),\ldots,I(N))\in S^N:
    \#\{t:I(t)=s\}=N/M\ \text{for every }s\in S\right\}.
\]
We call these edges \emph{candidates}.  They form a counting class;
we do not assume that the full joint type $\mathcal A_N$ can itself be
selected by local linear maps.

For a fixed candidate $I\in\mathcal A_N$, each column of a partner
$J\in S^N$ satisfying $J(t)-I(t)\in V$ can be chosen independently
from $C_V(I(t))$.  Hence the count over the \emph{entire original support}
is exactly
\begin{align}
  \#\bigl\{J\in S^N:J(t)-I(t)\in V\ \text{for all }1\leq t\leq N\bigr\}
  &=\prod_{s\in S}|C_V(s)|^{N/M}\notag\\
  &=\exp\bigl(NF_Z(V)\bigr)
   \leq\exp(Nr\beta),\qquad r=\dim V.
  \label{eq:fiber-count}
\end{align}
No joint-type or marginal-type condition is imposed on $J$.
In particular, \eqref{eq:fiber-count} also bounds the smaller class of
partners whose column differences span exactly $V$.

Finally, the size of the candidate class satisfies
\begin{equation}
  \frac{M^N}{(N+1)^M}
  \leq |\mathcal A_N|
  =\frac{N!}{((N/M)!)^M}
  \leq M^N.
  \label{eq:type-count}
\end{equation}
There are at most $(N+1)^M$ count vectors for words in an $M$-letter
alphabet.  The balanced vector maximizes its multinomial coefficient
when $M$ divides $N$: if two counts differ by at least two, transferring
one occurrence from the larger to the smaller strictly decreases the
factorial denominator.  At a maximum, all counts therefore equal
$N/M$.  Since the multinomial classes partition the $M^N$ words, their
largest class has size at least $M^N/(N+1)^M$.  The upper bound follows
from $\mathcal A_N\subseteq S^N$.
For fixed $\mathbf d$, equation~\eqref{eq:type-count} yields
\[
  \log|\mathcal A_N|=Nh-o(N).
\]
